\documentclass[a4paper,11pt]{scrartcl}
\usepackage[utf8x]{inputenc}
\usepackage[catalan]{babel}
\usepackage{titlesec}
% A ses llengües llatines, el primer paràgraf ha d'anar tabulat
\usepackage{indentfirst}
\usepackage{amsmath}
\usepackage{float}
\usepackage{graphicx}
\usepackage{subfigure}
\usepackage{booktabs}
\usepackage{multirow}
\usepackage{hyperref}
\usepackage{url}
\usepackage{multirow}
\usepackage{tabularx}

%% apt-get install python-pygments
%% wget http://minted.googlecode.com/hg/minted.sty
\usepackage{minted} %

%% aptitude install texlive-fonts-extra
\usepackage{newcent} %font mes wapa

\graphicspath{{diagrames/}}

% Estil de seccions
\titleformat{\section}{\large\sectfont}{\thesection}{1em}{}
\titleformat{\subsection}{\bfseries\sectfont}{\thesubsection}{1em}{}
% Estil numeracio subseccions http://help-csli.stanford.edu/tex/latex-sections.shtml#number
%\def\thesubsection{\alph{subsection})}

%% Tamany del codi font python inserit amb el minted
\newcommand{\codeSize}{\footnotesize}

\title{Pràctica de Cerca Local: \\ \huge{Carmageddon}}
\author{ Miquel Mas Fiol \thanks{mkelmas a gmail punt com} \\
	 Bartomeu Miró Mateu \thanks{bartomeumiro a gmail punt com} \\}

\begin{document}

  \maketitle

  \begin{abstract}
    \begin{center}
      \includegraphics[width=0.6\textwidth]{figures/carmageddon-logo.png}
    \end{center}
  \end{abstract}

  \newpage
  \setcounter{page}{2}
  \tableofcontents
  \newpage

  \section{Interpretació de l'enunciat}

En aquest problema tenim una ciutat, passatgers i conductors. Tots els passatgers i conductors
tenen un origen i un destí fixat dins la ciutat on han d'arribar i volem sabre quins conductors
transporten a quins passatgers, sabent la ruta que duen a terme.

Tots els passatgers poden ser recollits per un conductor que el portarà al seu destí,
així doncs un passatger no pot fer el seu tarjecte amb més de un conductor.

Hem de tenir amb compte que els conductors tenen una capacitat \emph{C} limitada del vehicle que no pot ser
sobrepassada. Això signfica que en un mateix instant no hi pot haver més de \emph{C} passatgers dins un vehicle,
no obstant durant tot el recorreut del conductor pot deixar i recollir diversos passatgers arribant a transportar
més de \emph{C} passatgers.


TODO DIAGRAMA

\emph{Tots} els conductors i passatgers han d'arribar \textbf{a temps} al seu destí. Donat que és
suposa una velocitat \emph{V} constant i que el conductor és l'últim arribar això significa aquest no podrà
superar més d'una determinada distància. En concret es presenta l'escenar de iniciar la ruta a les 7 del mati
i acabar a les 8. Això dona una hora re recorregut a velocitat V per el conductor. Si aquesta velocitat es
de 30km/h tenim que un conductor pot recorrer com a molt 30km en la seva ruta. Notem que en l'àmbit
d'aquest problema suposam que el cost temporal de pujar o baixar un passatger és despreciable.

Un conductor pot passar a ser passatger. En tal cas es comportarà com un passatger més. Dit cas
dependrà de la valoració que es faci d'eliminar un conductor o de si comparteix ruta amb un altre
conductor amb espai lliure i que per tant així es minimitza la distància recorreguda.

Donat aquest escenari tenim dos criteris de minimització.
Per una banda volem minimitzar la distància total recorreguda per tots els conductors i per l'altra
a més de minimitzar la distància també volem minimitzar el número de conductors.

Hem de tenir amb compte que dins el primer en certa mesura també s'inclou la minimització de conductors
ja que el còmput de distància es fa per número de conductors i en cas d'haver dos conductors compartint
ruta el fet d'eliminar-ne un resta la distància de la seva ruta minimitzant el còmput total.

En quan al segon criteri s'ha de preveure que número de conductors i kilòmetres són dues unitats de 
mesura diferents que han de donar lloc a una única valoració per tant s'haurà d'establir una manera 
de combinar-les o cercar-hi alguna relació.

\section{Paràmetres}

El número de persones total ve determinat per el paràmetre \emph{N}, de les quals \emph{M} són conductors i per
tant podem definir que \emph{P} és el número de passatgers. 

La ciutat és defineix com una matriu de 10x10 Km on cada illeta es de 100x100 m, així doncs
la ciutat és una matriu de 100x100 illetes.

La distància \emph{D} entre punts (\emph{o}rigen i \emph{d}estí) de la ciutat ve determinada per la funció:
$$ D(o,d) = |o_x - d_x| + |o_y - d_y| $$

Hem de tenir amb compte que la ciutat al ser una matriu per anar d'un punt A a un punt B hi ha múltiples
rutes amb la mateixa distància. Cosa que pot suposar un aventatge ja que així els únics paràmetres
d'una ruta son el punt d'origen, el destí i la distància.

  \newpage
  \section{Modelització del problema}
La pràctica ha estat desenvolupada en \emph{Python} el qual te una estructura de dades molt optimitzada
i integrada al llenguatge que és el diccionari\footnote{\url{http://docs.python.org/library/stdtypes.html#dict}}
el qual dona suport a la majoria d'estructura de dades de la pràctica.

\subsection{Representació de les dades}
Per la representació de les dades es fan servir 3 classes: passatger, conductor i estat.

On estat es l'estat requerit per les llibreries del AIMA. Passatger i conductor han estat
definits per facilitar l'elaboració de la pràctica.

\subsubsection{Passatger}
Cada passatger ve representat amb un identificador i els punts d'origen i destí tots ells són
camps invariables durant l'execució.

Assosciat a aquesta classe simplement hi ha els mètodes per obtenir el valors de dits camps.

\subsubsection{Conductor}
El conductor és una classe que hereda del passatger, així doncs també té el camp que l'identifica
el seu origen i destí. A més també té una capacitat màxima.
Per altra banda el conductor també inclou una llista dels noms dels passatgers que transportarà
durant tota la seva ruta.

Per gestionar aquesta classe són necessaris diveros mètodes més complicats que simples \emph{settes} i \emph{gettes}.
En primer lloc hi ha una funció \texttt{pickupPassenger} per recollir un passatger, aquesta afegeix el passatger
a la llista de passatgers a transportar. De manera complementària te \texttt{leavePassenger} per deixar un 
passatger.

Per tal de sabre quina és la ruta que fa el conductor hi ha la funció \texttt{getRoute} la qual
fa una ordenació dels punts clau (\texttt{checkpoints}) per on es passa i registra a cada punt si es
l'origen del condcutor, l'origen d'un passatger, el destí del passatger o el destí del condcutor. Així
doncs un ruta tindria l'aspecte seguent: 


TODO ruta


Associat a una ruta ens interessa sabre la distància que recorr, per això hi ha la funció \emph{getKm}
que torna tal valor.

En aquest punt s'ha de tenir amb compte com es generen les rutes ja que això infuleix directament en
si un passatger pot ser tranportat o no. La manera òptima en quan a recorregut seria fent un viatjant
de comerç entre tots els punts per on s'ha de passar tinguent amb compte dos factors:
Un cop s'agfa un passatger s'ha de introudir a la llista de punts el seu destí.
S'ha de vigilar de no agafar més passatgers del que pot transportar el vehicle
abans de deixar-ne algun.

Aquest problema tot hi ser TSP pot dur-se a terme ja que un conductor no sol transportar mes de 3
persones en una ruta la qual cosa suposen un total de 8 punts. Així i tot el càlcul de la ruta
es la funció mes emprada durant l'execució i interessa minimitzar el seu cost.

Per tal de simplificar-ho en lloc del viatjant de comerç es van cercant els punts més propers de
manera imediata.

TODO figura mes propers manera immediata VS viatjant de comperç

Donat també que hi ha poc punts i son molt dispersos aquesta aproximació, en el nostre cas,
dona un resultat acceptable estalviant temps d'execució.


\subsubsection{Ciutat}
La citutat en si no requereix representació, no té perquè estar guardara en memòria ja que senzillament
es pot guardar a cada passatger i conductor el seu punt d'inici i de destí com dues tuples \emph{XY}.

Per tant la so\lgem ució no és més que el conjunt de rutes dels conductors que formen part de l'estat.

\subsubsection{Estat}
L'estat ve determinat per un diccionari de conductors i un de passatgers. Ambdós diccionars estan
indexat per els noms de passatger i conductor respectivament però contenen coses diferents.

El diccionari de conductors conté un punter a cada objecte conductor mentre el el diccionari
de passatgers conté una tupla amb l'objecte del passatger i el nom del conductor que el transporta.

El fet de tenir un diccionari de passatgers es informació redundant però així es pot sabre
amb qui va un passatger en temps mig de  $O(1)$ en lloc de $O((n_{passatgers} + n_{conductors})/2)$ 
si haguéssim de recorrer tots els conductors i comprovar tots els passatgers que transporten.

TODO figura del dos diccionaris i aabaix tots els objectes

\subsection{Estat inicial}
Per la generació de l'estat inicial s'han contemplat dues maneres de fer-ho. En cap de les dues es contempla
la possiblitat de que algun passatger no sigui transportat per algun conductor. Com a invariant durant
tota l'execució tot passatger es transportat per algun conductor.

També en ambdós casos l'assignació d'un passatger a un conductor és aleatoria.

A continuació trobar les dues maneres d'inicialitzar el problema.

\subsubsection{Per saturació de conductors}
En aquesta inicialització del problema es recorren tots els passatgers assignant-los a conductors. No es passa
al següent conductor fins que aquest no arriba a un determinat nombre prefixat. Així doncs a un conductor no se
li assignara un passatger fins que tots els anteriors no siguin \emph{saturats}. Més endavant veurem
que el control de si el cotxe esta massa ple en algun instant es detectat per l'heurístic, així doncs
durant l'execució de la pràctica hi poden haver conductors que transportin més passatgers dels que
caben al vehicle.

Amb aquesta inicialització cap passatger queda sense transport, molt possiblement els condcutros escedeixin
la seva distància màxima i gairebé sempre hi haurà conductors que no transportin ningu.
Dit esces de la capacitat màxima es veura reflectit molt negativament a l'heurístic.

TODO diagrama 

\subsubsection{Evitant deixar conductors buits}
L'altre alternativa plantejada per la inicialització és recórrer tots els passatgers i assignar-los a conductors
de tal manera que no s'assigni un segon passatger a un conductor mentre hi hagin conductors buits.
D'aqueta manera s'afavoreix que les rutes dels conductors siguin més curtes i que hagi més
gent conduint.

TODO diagrama

\subsection{Operadors de transformació}
S'han establert quatre operadors de transformació agrupats en dos conjunts. Així doncs en la fase d'experimentació
es triara si l'agafa el conjunt 1 o el conjunt 2.

\subsubsection{Conjunt 1}
\subsubsection{Canvi de passatger de vehicle}

\subsubsection{Eliminació d'un conductor buit}

\subsubsection{Conjunt 2}
\subsubsection{Intercanvi de un passatger per un altre}
de vehicles diferents


\subsubsection{Eliminació d'un conductor no buit}

\subsection{Heurístic}
Es disposa de dos heurístics l'un per minimitzar el numero total de kilòmetres recorreguts i l'altre per
minimitzar el número de kilòmetres i a més també minimitzar el número de conductors.

En ambdós heurístics s'ha de tenir amb compte que hi pot haver estat que no compleixin les restriccions del problema
(com ara que no s'arriba a temps), en aquest cas l'heurístic ha de penalitzar suficient aquests casos com per no
ser tinguts amb compte.

  \newpage
  \section{Experimentació}

\subsection{Se\lgem ecció dels operadors}
Per a realitzar l'experiment 1 ha fixat en nombre d'usuaris a 200 i el nombre de conductors a 100.
Posteriorment s'ha executat 10 cops amb cada un dels conjunts d'operadors.


Els resultats són ja donats les mitjes són:

\begin{center}
\begin{tabular}{l|cccc}
         & Temps & Conductors & Distància & Estats (aprox)\\
\hline
Conjunt 1 & 0:02:14 & 98 & 1009304 & 9950 \\
Conjunt 2 & 0:01:30 & 98 & 1113629 & 5200
\end{tabular}
\end{center}


\begin{figure}[H]
\begin{center}
 \includegraphics[width=0.6\textwidth]{figures/test1-gr1.jpg}
 \label{test1-gr1}
\caption{Test 1}
\end{center}
\end{figure}
                

El conjunt d'operadors 1 produeix millor resultat que el conjunt 2.
Suposam que això és degut a que l'operador genera mes estats (gairebé el doble) per a cada iteració i simplement a més estats per triar s'agafen més bons.

Es pot veure que el conjunt 2 es notablement més ràpid però no en la mateix proporció que els estats generats. Això es deu a que aquests estats
son més costosos, entre d'altres coses perquè un intercanvi implica dues operacions mentre que un canvi només una, entre d'altres detalls d'ambdós operadors.


\subsection{Se\lgem ecció de la generació de l'estat inicial}
Per realitzar l'experiment 2 s'ha fixat el nombre d'usuaris i el nombre de conductors com a l'anterior experiment.                                                                                   
S'ha executat 10 cops amb cada estratègia de generació inicial.
                                                                                                                                                                                                     
Els resultats són:


\begin{center}
\begin{tabular}{l|ccc}
         & Temps & Conductors & Distància\\
\hline
\emph{all one first} & 0:01:51 & 98 & 974543  \\
\emph{full first} & 0:02:08 & 98 & 1030151 
\end{tabular}
\end{center}

\begin{figure}[H]
\begin{center}
 \includegraphics[width=0.6\textwidth]{figures/test2-gr1.jpg}
 \label{test2-gr1}
\caption{Test 2}
\end{center}
\end{figure}
                                                                                                                                                                                  
La estratègia de generació inicial \emph{all one first} produeix millors resultats tant en temps com distància,
aquest fet es deu a que te els passatgers més repartits en contraposició al \emph{full first} que els te
condensats i tarda més en repartir-los de manera adequada. De fet no te gaire sentit que els conductors
viatgin buits a l'inici ja que hem descartat el conjunt d'operadors que tal vegada haguessin pogut
jugar un bon paper amb el \emph{full first}.

\subsection{Se\lgem ecció dels paràmetres del Simulated Anealing}
Per a determinar els paràmetres que obtenen un millor resultat per a l'execució de l'algorisme \emph{simulated annealing},
hem executat reiterades vegades, aproximadament 10, l'algorisme amb diferents paràmetres per a $\lambda$ i \emph{k}, concretament s'ha fet amb
\texttt{[0.1,0.01,0.001,0.0001]} i \texttt{[1,5,25,125]} respectivament, fixant el nombre d'iteracions en 2000.

El millor resultat s'ha obtingut amb els paràmetres $k$ = 5 i $\lambda$ = 0.01, com indica el següent gràfic (fig. \ref{test3-gr1}).

\begin{figure}[H]
\begin{center} 
 \includegraphics[width=0.6\textwidth]{figures/test3-gr1.jpg}
\label{test3-gr1}
\caption{Evolució de $\lambda$}
\end{center}
\end{figure}

Fixats aquests dos paràmetres en els valors que han obtingut millor resultat, s'ha ajustat el nombre d'iteracions executant diverses 
vegades l'algorisme amb un nombre diferent d'iteracions. Es pot veure (fig. \ref{test3-gr2}) en el següent gràfic el millor valor per al nombre d'iteracions.

\begin{figure}[H]
\begin{center}
 \includegraphics[width=0.6\textwidth]{figures/test3-gr2.jpg}
 \label{test3-gr2}
\caption{Evolució de \emph{k}}
\end{center}
\end{figure}


Aleshores, podem comparar els resultats obtinguts amb simulated Annealing i Hill climbing (fig. \ref{test3-gr3} i \ref{test3-gr4}).


\begin{figure}[H]
\begin{center} 
 \includegraphics[width=0.6\textwidth]{figures/test3-gr3.jpg}
\label{test3-gr3}
\caption{}
\end{center}
\end{figure}


\begin{figure}[H]
\begin{center}
 \includegraphics[width=0.6\textwidth]{figures/test3-gr4.jpg}
 \label{test3-gr4}
\caption{}
\end{center}
\end{figure}

Es pot observar que el Simulated Annealing obté millors resultats que el Hill Climbing. Per contra, té un temps d'execució molt més elevat.


\subsection{Observació de l'evolució temporal del Hill Climbing}
Per a l'experiment 4 s'ha fixat el nombre d'usuaris i el nombre de conductors com a l'experiment inicial.
En aquest cas, s'utilitza el conjunt d'operadors 1, i l'algorisme de generació d'estats inicials \emph{all one first}.
S'ha executat 10 vegades cada simulació.

Es pot observar com el temps d'execució segueix clarament un patró quadràtic ascendent en funció del nombre d'usuaris
del servei.

\begin{figure}[H]
\begin{center}
 \includegraphics[width=0.6\textwidth]{figures/test4-gr1.jpg}
 \label{test4-gr1}
\caption{Evolució temporal en funció del nombre d'usuaris}
\end{center}
\end{figure}


\subsection{Comparació \emph{Hill Climbing} i \emph{Simulated Annealing} amb els dos heurístics}
Per a estimar la diferencia entre la distància total recorreguda per a \emph{Hill Climbing} i \emph{Simulated Annealing} per als dos heurístics,
hem realitzat 4 tests, cada un amb un heurístic i algorisme de cerca diferent, i l'hem executat 10 cops cada un.
Podem observar la taula amb els resultats.

\begin{center}
\begin{tabular}{l|llll}
 & \multicolumn{2}{c}{Hill Climbing} & \multicolumn{2}{c}{Simulated Annealing}\\
           & km       & veh       & km        & veh\\
\hline
temps      & 0:00:30  &  0:00:40  &  0:00:46  &  0:00:47\\
disntacia  & 549300   & 556103    &  857856   & 891356\\
conductors & 49       & 8         & 50        & 50\\
heurístic  & 549300   & 386556103 & 857856    & 500891356
\end{tabular}
\end{center}

\begin{figure}[H]
\begin{center}
 \includegraphics[width=0.6\textwidth]{figures/test6-gr1.jpg}
 \label{test6-gr1}
\caption{Compració HC vs SA}
\end{center}
\end{figure}

D'aquests resultats observam que l'algorisme \emph{Hill Climbing} troba millors solucions que el \emph{Simulated Annealing} si disposa d'un temps similar.
Suposam que augmentant el nombre d'iteracions del \emph{Simulated Annealing} podríem aconseguir millors resultats que amb el \emph{Hill climbing}.


S'observa també que els dos heurístics funcionen com esperàvem en el cas de \emph{Hill Climbing},(disminuint el nombre de km en un cas
i el nombre de km i conductors en l'altre), però en el \emph{Simulated Annealing} s'obté un resultat similar.
Es possible que això sigui degut a que els heurístics s'han ponderat utilitzat \emph{Hill Climbing}.

En aquest experiment sembla que ens decantaríem en favor de \emph{Hill Climbing}.






\subsection{Se\lgem ecció del valor de \emph{M} per millor inicialització del problema}
Per a realitzar l'experiment 7 s'ha se\lgem eccionat un problema de mida 50 usuaris, i s'ha provat per a totes les proporcions
(de tots conductors fins a un sol conductor), 10 cops cada una.

Es pot observar al següent gràfic que les so\lgem ucions més benevolents s'obtenen aproximadament per a 30 conductors (20 no conductors al gràfic),
que podem extrapolar aproximadament $\dfrac{3}{5}$ de conductors del total per a qualsevol problema.
Els valors en que no s'ha trobat so\lgem ució (a partir de 34) no es mostren.

\begin{figure}[H]
\begin{center}
 \includegraphics[width=0.6\textwidth]{figures/test7-gr1.jpg}
 \label{test7-gr2}
\end{center}
\caption{Evolució del la distància}
\end{figure}

Per al temps de trobar una so\lgem ució, podem observar que on més temps es tarda es a una proporció de $\dfrac{2}{5}$ conductors.
Intuïm que es la zona on es generen mes intercanvis de passatgers, ja que n'hi ha suficients per a generar molts d'estats,
però no massa com per a que molts d'estats no siguin vàlids.
Amb una proporció més alta de passatgers es triga menys temps, però com hem explicat anteriorment es difícil trobar una so\lgem ució que satisfaci les restriccions.

\begin{figure}[H]
\begin{center}
 \includegraphics[width=0.6\textwidth]{figures/test7-gr2.jpg}
 \label{test7-gr2}
\end{center}
\caption{Evolució del temps}
\end{figure}
        

  \newpage
  \section{Manual d'usuari}

La practica està escrita en \emph{python}\footnotemark, per tant és necessari tenir
l'interpret insta\lgem at.
Es recomana l'ús de \emph{pypy}\footnotemark en lloc de l'interpret per defecte ja 
que aquest implementa un \emph{just-in-time complier} i accelera notablement l'execució
de la pràctica.

Per dur a terme l'execució s'ha de fer amb el fitxer \texttt{prog.py} si es executat
sense parametres mostrarà quins suporta i alguns exemples d'execució.

Essencialment te dues maneres de funcionar, una es executar un test dels demanats a l'enunciat,
l'altre es executar un hill climbing o simulated anealing amb determinats paràmetres.

\subsection{Execució d'un test}
Per executar un test simplement s'ha de passar com a paràmetre el número del test a executar i 
aquest s'executara i mostrarà per pantalla el resultat.

En cas que sigui un test amb diferents opcions per defecte s'executa amb la opció que dona
una millor execució. Per tal de canviar dites opcions s'ha d'editar el fitxer prog.py anant
a la capçalera de la funció i canviant el paràmetre per defecte.

\begin{minted}[fontsize=\small]{python}
def test4_TemporalEvolution(operatorSet="1", initDistrib="allOneFirst"):
   ...
\end{minted}

\subsection{Execució simple}
Si sols volem executar una una cerca simple tenim dues opcions, una es carregar un fitxer de configuració
guardat i l'altre espeficicar tots els paràmetres.

Els paràmetres per a fer una execució simple són:

\begin{minted}[fontsize=\small]{bash}
pypy prog.py N M inicialitzacio heuristic conjuntOperadors algorismeCerca

\small
\begin{center}
\begin{tabularx}{\textwidth}{llX}
\textbf{parametre}  & \textbf{Valors} & \textbf{Descripció}\\
\midrule
N  & valor enter & Numero d'usuaris del servei\\
M  & valor enter & NUmero d'usuaris no conductors\\
\midrule
inicialització & allOneFirst & REF\\
               & fullFirst   & REF\\
\midrule
heuristic      & km  & Minimitza el número de km recorreguts\\
	       & veh & Mínimitza el numero de km i el numero de conductors\\
\midrule
conjuntOperadors      & 1 & Moure passatgers i llevar conductors buits\\
		      & 2 & Intercanviar passatgers i llevar conductors\\
\midrule
algorismeCerca & hillClimbing & Executa un Hill Climbing\\
               & simulatedAnnealing & Executa Simulated Annealing
\end{tabularx}
\end{center}
\normalsize


Els paràmetres per fer una execució des de un fitxer són:

\begin{minted}[fontsize=\small]{bash}
pypy prog.py fitxer.cfg heursistic conjuntOperadors algorismeCerca
\end{minted}

On fitxer es el nom del \texttt{fitxer.cfg} de configuració i la resta de paràmetres són
iguals a l'apartat anterior.

\footnotetext{\url{http://python.org}}
\footnotetext{El lloc oficial és \url{http://pypy.org/} però es pot
trobar precompilat i empaquetat per Debian/Ubuntu a \url{https://launchpad.net/~pypy/+archive/ppa}}


  \newpage
  \input{peu}
  \newpage

  \section{Annex: Codi font}

\subsection{Carmageddon}
\inputminted[linenos, frame=lines, fontsize=\codeSize]{python}{../aima/carmageddon.py}
\newpage

\subsection{State}
\inputminted[linenos, frame=lines, fontsize=\codeSize]{python}{../aima/state.py}
\newpage

\subsection{Driver}
\inputminted[linenos, frame=lines, fontsize=\codeSize]{python}{../aima/driver.py}
\newpage

\subsection{Passenger}
\inputminted[linenos, frame=lines, fontsize=\codeSize]{python}{../aima/passenger.py}
\newpage


\subsection{showLog}
\inputminted[linenos, frame=lines, fontsize=\codeSize]{python}{../aima/showLog.py}


\subsection{Prog}
\inputminted[linenos, frame=lines, fontsize=\codeSize]{python}{../aima/prog.py}



\end{document}
